\begin{table}[htbp]
\centering
\small
\caption{Problema G (openai/gpt-4o-mini): Fase E, acurácia do LLM contra o rótulo real por número de exemplos.}
\label{tab:bloco3_G_faseE_nshot_gpt4mini}
\begin{tabular}{lllrrrrr}
\toprule
Subproblema & Classes & Atributos & n = 0 & n = 4 & n = 10 & n = 20 & n = 40 \\
\midrule
G1 (normalizado) & A / B & peso/altura & 25,6 (22,3) & 19,6 (24,1) & 16,7 (25,0) & 12,2 (26,3) & 13,3 (26,0) \\
G1 (normalizado) & A / B & x1/x2 & 30,0 (21,0) & 20,7 (23,8) & 21,1 (23,7) & 16,3 (25,1) & 15,9 (25,2) \\
\textbf{G2 (kg e m)} & \textbf{Mulher / Homem} & \textbf{peso\_kg/altura\_m} & \textbf{88,9 (3,3)} & \textbf{83,7 (4,9)} & \textbf{82,2 (5,3)} & \textbf{81,5 (5,6)} & \textbf{78,9 (6,3)} \\
G2 (kg e m) & Mulher / Homem & x1/x2 & 62,2 (11,3) & 62,2 (11,3) & 69,3 (9,2) & 71,1 (8,7) & 68,5 (9,4) \\
G2 (kg e m) & A / B & peso\_kg/altura\_m & 40,0 (18,0) & 51,5 (14,6) & 48,9 (15,3) & 43,7 (16,9) & 40,4 (17,9) \\
\bottomrule
\multicolumn{8}{p{0.95\linewidth}}{\footnotesize Fase E nas 30 amostras de teste de cada semente; o especialista é a melhor métrica da Fase A (eleita por acurácia vs real). Acurácia do LLM contra o rótulo real, erros em 30 entre parênteses; médias das sementes e sorteios. Em negrito: G2 com nomes semânticos.} \\
\end{tabular}
\end{table}
